\chapter{Timeline}

The timeline is displayed at the top of each waveform group and shows the X axis scale for the group. The timeline (and
all accompanying waveform views in the group) may be zoomed by scrolling with the mouse wheel, or panned by dragging
with the left mouse button. Pressing the middle mouse button (or double-clicking or double-tapping) auto-scales the
group so that the entire waveform is visible.

With the fixed divisions grid (see \hyperref[sec:grid]{Grid}), the timeline is split into a fixed number of
divisions. The center division is labeled with its time and the others with their offset from it (e.g. ``+200 ns"),
and the mouse wheel steps the time per division through the 1-2-5 sequence. Auto-scaling rounds the time per division
up to the next value in the sequence and centers the waveform, leaving some space on either side of it.

Unlike classical oscilloscope user interfaces, there is \emph{no relationship} between the timeline scale or position
and the duration of the acquisition. It is possible to zoom or scroll beyond the end of the acquisition (displaying
empty background with no signal) or have a deep capture in which nearly all acquired data is offscreen.

Note that the timeline may occasionally show units other than time. For example, an ``eye width" measurement has X axis
units of voltage and Y axis units of time, and a spectrum analyzer channel has X axis units of frequency.


\begin{figure}[h]
\centering
\includegraphics[width=14cm]{images/timeline.png}
\caption{The timeline}
\label{timeline}
\end{figure}

The position of the trigger event is marked by a downward-pointing arrow on the timeline, color coded to match the
channel selected as the primary trigger source. The trigger arrow cannot be interacted with currently, but in the future
(scopehal-apps:173) it will be draggable to adjust the trigger position.
